%\makeglossaries
\makenoidxglossaries

% Danh mục thuật ngữ và từ viết tắt
\newglossaryentry{API}{
    type=\acronymtype,
    name={API},
    description={Giao diện lập trình ứng dụng (Application Programming Interface)},
    first={API}
}

\newglossaryentry{BFF}{
    type=\acronymtype,
    name={BFF},
    description={Kiến trúc backend phục vụ riêng cho frontend (Backend For Frontend)},
    first={Backend For Frontend (BFF)}
}

\newglossaryentry{CDN}{
    type=\acronymtype,
    name={CDN},
    description={Mạng phân phối nội dung (Content Delivery Network)},
    first={Content Delivery Network (CDN)}
}

\newglossaryentry{CICD}{
    type=\acronymtype,
    name={CI/CD},
    description={Tích hợp liên tục và triển khai liên tục (Continuous Integration / Continuous Deployment)},
    first={Continuous Integration / Continuous Deployment (CI/CD)}
}

\newglossaryentry{CORS}{
    type=\acronymtype,
    name={CORS},
    description={Cơ chế chia sẻ tài nguyên giữa các nguồn gốc khác nhau (Cross-Origin Resource Sharing)},
    first={Cross-Origin Resource Sharing (CORS)}
}

\newglossaryentry{CRUD}{
    type=\acronymtype,
    name={CRUD},
    description={Bốn thao tác cơ bản trên dữ liệu: Tạo, Đọc, Cập nhật, Xoá (Create, Read, Update, Delete)},
    first={Create, Read, Update, Delete (CRUD)}
}

\newglossaryentry{DTO}{
    type=\acronymtype,
    name={DTO},
    description={Đối tượng truyền dữ liệu giữa các tầng (Data Transfer Object)},
    first={Data Transfer Object (DTO)}
}

\newglossaryentry{GV}{
    type=\acronymtype,
    name={GV},
    description={Giáo viên},
    first={Giáo viên (GV)}
}

\newglossaryentry{GVCN}{
    type=\acronymtype,
    name={GVCN},
    description={Giáo viên chủ nhiệm},
    first={Giáo viên chủ nhiệm (GVCN)}
}

\newglossaryentry{GVBM}{
    type=\acronymtype,
    name={GVBM},
    description={Giáo viên bộ môn},
    first={Giáo viên bộ môn (GVBM)}
}

\newglossaryentry{HTTP}{
    type=\acronymtype,
    name={HTTP},
    description={Giao thức truyền siêu văn bản (HyperText Transfer Protocol)},
    first={HyperText Transfer Protocol (HTTP)}
}

\newglossaryentry{HTTPS}{
    type=\acronymtype,
    name={HTTPS},
    description={Giao thức truyền siêu văn bản bảo mật (HyperText Transfer Protocol Secure)},
    first={HyperText Transfer Protocol Secure (HTTPS)}
}

\newglossaryentry{JPA}{
    type=\acronymtype,
    name={JPA},
    description={Đặc tả giao diện lập trình Java cho ánh xạ quan hệ đối tượng (Java Persistence API)},
    first={Java Persistence API (JPA)}
}

\newglossaryentry{JWT}{
    type=\acronymtype,
    name={JWT},
    description={Chuẩn mã hoá token xác thực phân tán (JSON Web Token)},
    first={JSON Web Token (JWT)}
}

\newglossaryentry{ORM}{
    type=\acronymtype,
    name={ORM},
    description={Kỹ thuật ánh xạ đối tượng sang bảng quan hệ (Object-Relational Mapping)},
    first={Object-Relational Mapping (ORM)}
}

\newglossaryentry{REST}{
    type=\acronymtype,
    name={REST},
    description={Phong cách kiến trúc dịch vụ web phi trạng thái (Representational State Transfer)},
    first={Representational State Transfer (REST)}
}

\newglossaryentry{SLA}{
    type=\acronymtype,
    name={SLA},
    description={Thoả thuận mức độ dịch vụ (Service Level Agreement)},
    first={Service Level Agreement (SLA)}
}

\newglossaryentry{SSR}{
    type=\acronymtype,
    name={SSR},
    description={Kỹ thuật render trang web phía máy chủ (Server-Side Rendering)},
    first={Server-Side Rendering (SSR)}
}

\newglossaryentry{TKB}{
    type=\acronymtype,
    name={TKB},
    description={Thời khoá biểu},
    first={Thời khoá biểu (TKB)}
}

\newglossaryentry{TLS}{
    type=\acronymtype,
    name={TLS},
    description={Giao thức bảo mật tầng truyền tải (Transport Layer Security)},
    first={Transport Layer Security (TLS)}
}

\newglossaryentry{UC}{
    type=\acronymtype,
    name={UC},
    description={Ca sử dụng},
    first={Ca sử dụng (UC)}
}

\newglossaryentry{UML}{
    type=\acronymtype,
    name={UML},
    description={Ngôn ngữ mô hình hoá thống nhất (Unified Modeling Language)},
    first={Unified Modeling Language (UML)}
}

\newglossaryentry{XSS}{
    type=\acronymtype,
    name={XSS},
    description={Tấn công chèn mã độc phía trình duyệt (Cross-Site Scripting)},
    first={Cross-Site Scripting (XSS)}
}
